\scanpage{90}

% L1–4
est un torseur sous $G$. Le choix d'un $r\in R$ définit un \deleted{\unsure{épim|épim.}} $G_n\xrightarrow{\varphi_r}G$, d'où un système \inserted{$\unsure{\tau_*|\mathrm T_*}$} $\underline\tau_*^r=(\underline\tau_i^r)=\varphi_r(\underline\sigma_i))$, transformé en un syst. conjugué \inserted{par $\operatorname{int}(g)$} si on change $r$ en $r'=gr$ ($g\in G$). \unsure{Si|\deleted{Si}} on se donne un hom

% L5
\[
\varphi_G:G\longrightarrow\operatorname{Aut}X,
\]

% L6–8
dire que $\tau_*^r=(\tau_i^r)=\varphi_G(\underline\tau_i^r)=(\varphi_G(\underline\tau_i^r))$ est formé de \uline{quasi-réflexions}, et que la condition de non dég. de B est satisfaite, ne dépend pas du choix de $r$.

% L9–11
De $=$ si on a une repr. \unsure{régulière} $(\varphi_G,\inserted{\varphi_*=}(\varphi_i)_{0\leq i\leq n}\quad(\varphi_i:\Phi_i(\Pi)\longrightarrow\operatorname{Grass}_i(X)))$ de $\Pi$, d'où $r=\varphi_*(\underline r)=X_*$, dire que $(\tau_{\unsure{j|*}}^r,X_*))$ \deleted{est} est non dég ne dépend pas du choix de $\underline r$ -- \deleted{satisfaire} \inserted{d'ailleurs}

% L12–16
les conditions \inserted{a) a$'$) b)} d) de A) sont intrinsèques. La première est \unsure{donc|en vertu} néc. et suff. pour l'existence loc$^{\mathrm{le}}$ (Zar) d'une \textquotedblleft réal. géom non dég\textquotedblright\ de $\Pi$ associée à $\varphi_G$. \deleted{\unsure{De même,|Donc} Ici encore, on voit que le groupe} \inserted{Comme $G$ est \uline{simplem$^{\mathrm t}$} \uline{transitif} sur} $R$, on voit d'ailleurs que tt objet géom. qu'on

% L17–20
peut associer intrinsèquement à une situation $\tau_*$ resp $(\tau_*,X_*)$ non dég [ tel que $D_0$, et ses sections $c,d$ ] est défini intrinsèquement aussi en termes d'une \inserted{repr.} $\varphi_G$ \inserted{quasi-polyédrale n.d.} resp. d'une \deleted{\unsure{$\varphi$}} réalisation géom \inserted{$(\varphi_G,\varphi_*)$} de $\Pi$.

% L21–22
Ceci dit, $\operatorname{Aut}\varphi_G$ opère sur $(D_0,c,d)$, et on trouve

% L22
\[
\underline{\operatorname{Aut}}(X,\varphi_G)\xrightarrow{\sim}
\underline{\operatorname{Aut}}\bigl(D_0;\overset{\text{id sur}}{(c)\cdot(d)}\bigr)
\]

% L23–25
donc la cat. des $\varphi_G$ proj.le non dég. équivaut à celle des couples $\deleted{\unsure{(\varphi_G}}\bigl(\underbrace{(\varphi_G,\varphi_*)}_{\substack{\text{réal. géom proj.}\\\text{n. dég. de }\Pi}},T\bigr)$, $T$ un torseur sous $\underline{\operatorname{Aut}}(D_0,c,d)$ -- ou la \unsure{$=$|même} interprétation

% archive sheet number: 89
